

\subsection{Controlled ablation on the CMU20 benchmark}

We evaluated the agent on the completed 20-case CMU benchmark, spanning H, NNH, OH, CH$_2$CH$_2$OH, OCHCH$_3$, and ONN(CH$_3$)$_2$ adsorption across increasingly difficult intermetallic and monometallic surfaces.
For each case, we executed five runtime variants under a fixed MACE-small CPU backend: \textit{Full}, \textit{$-$Slip}, \textit{$-$Forbid}, \textit{$-$Term}, and \textit{1-Shot}.
Here \textit{$-$Slip} removes slip-conditioned guidance together with its associated FORBID constraint, \textit{$-$Forbid} keeps slip detection and history but removes the explicit ``do not retry this site'' instruction, \textit{$-$Term} disables convergence-based early stopping, and \textit{1-Shot} reduces the total budget to a single attempt.
The complete 20$\times$5 matrix was executed independently on four LLM backends---Gemini~2.5~Pro, Grok-4, GPT-5.4, and Claude~Sonnet~4.6---yielding 393 successful runs out of 400 attempts.

\begin{table}[t]
\centering
\scriptsize
\caption{Cross-backend ablation summary on the completed 20-case CMU benchmark. Each row aggregates one variant across four LLM backends. $\Delta E = E_{\mathrm{variant}} - E_{\mathrm{full}}$; positive values indicate worse adsorption energy than the same-backend full AdsMind agent. For 1-Shot, energy statistics are conditional on successful non-dissociated runs; failed attempts are retained in the success-rate denominator.}
\label{tab:ablation-summary-main}
\resizebox{\textwidth}{!}{%
\begin{tabular}{lrrrrrr}
\hline
Variant & Attempts & Successful & Mean $\Delta E$ vs full (eV) & Median $\Delta E$ (eV) & Mean iterations & Success rate \\
\hline
Full & 80 & 80 & +0.000 & 0.000 & 4.18 & 100.0\%  \\
$-$Slip & 80 & 80 & $-$0.008 & 0.000 & 4.33 & 100.0\% \\
$-$Forbid & 80 & 80 & $-$0.007 & 0.000 & 4.30 & 100.0\% \\
$-$Term & 80 & 80 & $-$0.036 & 0.000 & 4.58 & 100.0\% \\
1-Shot & 80 & 73 & +0.217 & 0.125 & 1.00 & 91.2\% \\
\hline
\end{tabular}%
}
\end{table}

\paragraph{CMU20 completion and natural 1-Shot failures.}
The newly completed cases (03, 06, 07, 08, 11) close the CMU benchmark gap across all five variants, not only the Full variant.
The Full, $-$Slip, $-$Forbid, and $-$Term variants contain 320/320 successful runs across CMU20.
The only failures occur in the 1-Shot control: case~06 fails on all four backends and case~08 fails on Grok-4, GPT-5.4, and Claude, yielding 73/80 successful 1-Shot runs.
These failures are not external errors; in each case the relaxation/analyzer completes, but the final structure is rejected because \texttt{is\_dissociated=True} with \texttt{dissociation\_count=1} and no valid best adsorption energy.
Across CMU20, the mean four-backend range is 0.153~eV, the median range is 0.029~eV, 8/20 cases agree within 0.01~eV, and 12/20 agree within 0.05~eV.
The largest disagreement occurs on case~08 (Ru$_3$Mo + NNH), where GPT-5.4 reaches $-7.191$~eV, Claude reaches $-6.869$~eV, and Gemini/Grok-4 remain at $-6.414$~eV.
Thus, the Full agent is reliable on CMU20, while backend agreement remains case-dependent rather than universal.

\begin{table}[t]
\centering
\caption{Full-row backend robustness within the completed CMU20 five-variant ablation matrix.}
\label{tab:cmu20-full-completion}
\begin{tabular}{lr}
\toprule
Metric & Value \\
\midrule
Successful Full runs & 80/80 \\
Mean four-backend range & 0.153~eV \\
Median four-backend range & 0.029~eV \\
Maximum four-backend range & 0.776~eV (case~08) \\
Cases within 0.01~eV & 8/20 \\
Cases within 0.05~eV & 12/20 \\
\bottomrule
\end{tabular}
\end{table}

The detailed per-backend energy matrices and formal pairwise tests below preserve the original locked 15-case CMU subset used in the first complete-energy ablation package.
We do not recompute ordinary 20-case Wilcoxon tests over the CMU20 1-Shot row because seven 1-Shot attempts are censored by natural dissociation/reaction failures; the completed CMU20 matrix is therefore summarized through success rates, failure audit, and success-conditioned energy deltas in the Supporting Information.

% ===== inlined from research/results/paper_tables.tex =====
% ============================================================
% AdsMind Paper-Ready Tables (15-case CMU benchmark)
% Regenerated from 4-backend × 15-case ablation matrices
% ============================================================

% SI Table: Gemini 15-case ablation energy matrix
\begin{table}[h]
\centering
\small
\caption{Gemini~2.5~Pro ablation on the 15 locked CMU benchmark cases. Lower (more negative) adsorption energy is better; bold marks the best variant per case.}
\label{tab:si-ablation-gemini}
\begin{tabular}{lrrrrr}
\toprule
Case & Full & $-$Slip & $-$Forbid & $-$Term & 1-Shot \\
\midrule
01 & \textbf{$-$3.632} & \textbf{$-$3.632} & \textbf{$-$3.632} & \textbf{$-$3.632} & $-$3.631 \\
02 & \textbf{$-$4.766} & \textbf{$-$4.766} & \textbf{$-$4.766} & \textbf{$-$4.766} & \textbf{$-$4.766} \\
04 & $-$2.060 & \textbf{$-$2.195} & $-$2.039 & $-$1.845 & $-$1.658 \\
05 & \textbf{$-$2.962} & \textbf{$-$2.962} & \textbf{$-$2.962} & \textbf{$-$2.962} & $-$2.947 \\
09 & \textbf{$-$1.974} & \textbf{$-$1.974} & \textbf{$-$1.974} & \textbf{$-$1.974} & $-$1.801 \\
10 & \textbf{$-$2.718} & \textbf{$-$2.718} & \textbf{$-$2.718} & \textbf{$-$2.718} & $-$2.122 \\
12 & \textbf{$-$2.351} & \textbf{$-$2.351} & \textbf{$-$2.351} & \textbf{$-$2.351} & $-$1.607 \\
13 & \textbf{$-$2.825} & \textbf{$-$2.825} & \textbf{$-$2.825} & \textbf{$-$2.825} & $-$2.376 \\
14 & \textbf{$-$3.617} & \textbf{$-$3.617} & \textbf{$-$3.617} & \textbf{$-$3.617} & \textbf{$-$3.617} \\
15 & \textbf{$-$3.258} & \textbf{$-$3.258} & $-$2.964 & \textbf{$-$3.258} & $-$2.889 \\
16 & \textbf{$-$4.663} & \textbf{$-$4.663} & \textbf{$-$4.663} & \textbf{$-$4.663} & $-$4.538 \\
17 & \textbf{$-$3.222} & $-$3.064 & $-$3.064 & $-$3.064 & $-$2.773 \\
18 & $-$2.779 & $-$2.779 & \textbf{$-$3.129} & \textbf{$-$3.129} & $-$2.418 \\
19 & $-$3.841 & $-$3.841 & \textbf{$-$4.042} & $-$3.594 & $-$3.133 \\
20 & $-$11.652 & \textbf{$-$11.830} & $-$11.652 & $-$11.652 & $-$10.604 \\
\bottomrule
\end{tabular}
\end{table}

% SI Table: Grok-4 15-case ablation energy matrix
\begin{table}[h]
\centering
\small
\caption{Grok-4 ablation on the 15 locked CMU benchmark cases. Lower (more negative) adsorption energy is better; bold marks the best variant per case.}
\label{tab:si-ablation-grok}
\begin{tabular}{lrrrrr}
\toprule
Case & Full & $-$Slip & $-$Forbid & $-$Term & 1-Shot \\
\midrule
01 & \textbf{$-$3.632} & \textbf{$-$3.632} & \textbf{$-$3.632} & \textbf{$-$3.632} & $-$3.631 \\
02 & \textbf{$-$4.766} & \textbf{$-$4.766} & $-$4.760 & \textbf{$-$4.766} & $-$4.150 \\
04 & \textbf{$-$2.255} & $-$2.039 & $-$1.845 & \textbf{$-$2.255} & $-$1.658 \\
05 & \textbf{$-$2.962} & \textbf{$-$2.962} & \textbf{$-$2.962} & \textbf{$-$2.962} & $-$2.947 \\
09 & \textbf{$-$1.974} & \textbf{$-$1.974} & \textbf{$-$1.974} & \textbf{$-$1.974} & \textbf{$-$1.974} \\
10 & \textbf{$-$2.718} & \textbf{$-$2.718} & \textbf{$-$2.718} & \textbf{$-$2.718} & $-$2.712 \\
12 & \textbf{$-$2.351} & \textbf{$-$2.351} & \textbf{$-$2.351} & \textbf{$-$2.351} & $-$2.081 \\
13 & \textbf{$-$2.825} & \textbf{$-$2.825} & \textbf{$-$2.825} & \textbf{$-$2.825} & $-$2.376 \\
14 & \textbf{$-$3.617} & \textbf{$-$3.617} & \textbf{$-$3.617} & \textbf{$-$3.617} & $-$3.245 \\
15 & \textbf{$-$3.258} & \textbf{$-$3.258} & \textbf{$-$3.258} & \textbf{$-$3.258} & $-$2.587 \\
16 & \textbf{$-$4.663} & \textbf{$-$4.663} & \textbf{$-$4.663} & \textbf{$-$4.663} & $-$4.552 \\
17 & $-$2.866 & \textbf{$-$3.146} & $-$3.064 & $-$3.064 & $-$3.066 \\
18 & \textbf{$-$3.129} & $-$2.779 & \textbf{$-$3.129} & \textbf{$-$3.129} & $-$2.418 \\
19 & \textbf{$-$4.045} & $-$3.594 & $-$3.594 & $-$4.027 & $-$3.594 \\
20 & \textbf{$-$11.653} & \textbf{$-$11.653} & $-$11.652 & $-$11.652 & $-$10.603 \\
\bottomrule
\end{tabular}
\end{table}

% SI Table: GPT-5.4 15-case ablation energy matrix
\begin{table}[h]
\centering
\small
\caption{GPT-5.4 ablation on the 15 locked CMU benchmark cases. Lower (more negative) adsorption energy is better; bold marks the best variant per case.}
\label{tab:si-ablation-gpt54}
\begin{tabular}{lrrrrr}
\toprule
Case & Full & $-$Slip & $-$Forbid & $-$Term & 1-Shot \\
\midrule
01 & $-$3.627 & $-$3.627 & $-$3.627 & \textbf{$-$3.630} & $-$3.557 \\
02 & \textbf{$-$4.769} & $-$4.764 & \textbf{$-$4.769} & \textbf{$-$4.769} & $-$4.343 \\
04 & \textbf{$-$2.255} & \textbf{$-$2.255} & \textbf{$-$2.255} & \textbf{$-$2.255} & $-$2.243 \\
05 & \textbf{$-$2.962} & \textbf{$-$2.962} & \textbf{$-$2.962} & \textbf{$-$2.962} & $-$2.843 \\
09 & \textbf{$-$1.991} & \textbf{$-$1.991} & \textbf{$-$1.991} & \textbf{$-$1.991} & \textbf{$-$1.991} \\
10 & $-$2.710 & $-$2.710 & $-$2.710 & \textbf{$-$2.718} & \textbf{$-$2.718} \\
12 & $-$2.081 & $-$2.081 & \textbf{$-$2.351} & \textbf{$-$2.351} & $-$2.081 \\
13 & \textbf{$-$2.825} & \textbf{$-$2.825} & \textbf{$-$2.825} & \textbf{$-$2.825} & $-$2.808 \\
14 & \textbf{$-$3.591} & $-$3.431 & \textbf{$-$3.591} & \textbf{$-$3.591} & $-$3.431 \\
15 & $-$3.258 & $-$2.964 & \textbf{$-$3.385} & $-$3.258 & $-$2.964 \\
16 & $-$4.272 & \textbf{$-$4.663} & \textbf{$-$4.663} & \textbf{$-$4.663} & \textbf{$-$4.663} \\
17 & \textbf{$-$3.241} & \textbf{$-$3.241} & \textbf{$-$3.241} & \textbf{$-$3.241} & $-$2.692 \\
18 & $-$2.724 & \textbf{$-$3.215} & \textbf{$-$3.215} & \textbf{$-$3.215} & $-$2.519 \\
19 & $-$4.013 & $-$4.100 & \textbf{$-$4.149} & $-$4.013 & $-$4.013 \\
20 & $-$11.652 & \textbf{$-$11.653} & $-$11.650 & $-$11.652 & $-$11.652 \\
\bottomrule
\end{tabular}
\end{table}

% SI Table: Claude 15-case ablation energy matrix
\begin{table}[h]
\centering
\small
\caption{Claude~Sonnet~4.6 ablation on the 15 locked CMU benchmark cases. Lower (more negative) adsorption energy is better; bold marks the best variant per case.}
\label{tab:si-ablation-claude}
\begin{tabular}{lrrrrr}
\toprule
Case & Full & $-$Slip & $-$Forbid & $-$Term & 1-Shot \\
\midrule
01 & \textbf{$-$3.630} & \textbf{$-$3.630} & $-$3.624 & \textbf{$-$3.630} & $-$3.557 \\
02 & \textbf{$-$4.769} & \textbf{$-$4.769} & \textbf{$-$4.769} & \textbf{$-$4.769} & $-$4.343 \\
04 & \textbf{$-$2.285} & $-$2.283 & \textbf{$-$2.285} & $-$2.243 & $-$1.845 \\
05 & \textbf{$-$2.962} & \textbf{$-$2.962} & \textbf{$-$2.962} & \textbf{$-$2.962} & \textbf{$-$2.962} \\
09 & \textbf{$-$1.991} & \textbf{$-$1.991} & \textbf{$-$1.991} & \textbf{$-$1.991} & $-$1.804 \\
10 & \textbf{$-$2.718} & \textbf{$-$2.718} & \textbf{$-$2.718} & \textbf{$-$2.718} & $-$2.710 \\
12 & \textbf{$-$2.351} & \textbf{$-$2.351} & \textbf{$-$2.351} & \textbf{$-$2.351} & $-$2.081 \\
13 & $-$2.825 & $-$2.825 & $-$2.825 & \textbf{$-$2.830} & $-$2.825 \\
14 & $-$3.580 & \textbf{$-$3.591} & \textbf{$-$3.591} & \textbf{$-$3.591} & $-$3.453 \\
15 & \textbf{$-$3.258} & \textbf{$-$3.258} & \textbf{$-$3.258} & \textbf{$-$3.258} & $-$2.767 \\
16 & \textbf{$-$4.663} & $-$4.272 & \textbf{$-$4.663} & \textbf{$-$4.663} & \textbf{$-$4.663} \\
17 & \textbf{$-$3.241} & \textbf{$-$3.241} & \textbf{$-$3.241} & \textbf{$-$3.241} & $-$2.692 \\
18 & $-$2.724 & \textbf{$-$3.215} & $-$2.838 & \textbf{$-$3.215} & $-$2.519 \\
19 & $-$4.013 & \textbf{$-$4.149} & \textbf{$-$4.149} & $-$4.013 & $-$4.013 \\
20 & \textbf{$-$11.653} & $-$11.652 & $-$11.652 & $-$11.652 & $-$11.652 \\
\bottomrule
\end{tabular}
\end{table}

% Table: Cross-LLM robustness (15-case, 4-backend)
\begin{table}[t]
\centering
\caption{Cross-LLM robustness across the 15-case ablation matrix. The four-backend range is max--min energy across Gemini, Grok-4, GPT-5.4, and Claude for the same case and runtime variant.}
\label{tab:cross-llm-summary}
\begin{tabular}{lrrr}
\toprule
Variant & Mean 4-backend range (eV) & $\leq$0.01\,eV & Largest range (eV) \\
\midrule
Full & 0.129 & 7/15 & 0.405 \\
$-$Slip & 0.184 & 5/15 & 0.555 \\
$-$Forbid & 0.136 & 8/15 & 0.555 \\
$-$Term & 0.077 & 9/15 & 0.433 \\
1-Shot & 0.425 & 0/15 & 1.050 \\
\midrule
Overall & 0.190 & 29/75 & 1.050 \\
\bottomrule
\end{tabular}
\end{table}

% Table: Key evaluation metrics (15-case)
\begin{table}[t]
\centering
\caption{Key evaluation metrics across all four backends on the 15-case CMU benchmark.}
\label{tab:key-evaluation-metrics}
\begin{tabular}{p{2.8cm}p{1.3cm}p{1.8cm}p{1.8cm}p{1.8cm}p{1.8cm}p{1.0cm}}
\toprule
Metric & Target & Gemini & Grok-4 & GPT-5.4 & Claude & Met \\
\midrule
Friedman $p$ & $< 0.05$ & 1.40e-06 & 2.21e-06 & 2.15e-03 & 1.14e-04 & 4/4 \\
H$_1$ support & $\geq$75\% & 7/11 (64\%) & 8/11 (73\%) & 7/11 (64\%) & 7/11 (64\%) & --- \\
BH $p$ (1-Shot) & $< 0.05$ & 0.0059 & 0.0115 & 0.0840 & 0.0134 & 3/4 \\
Backend conv. & Lower & \multicolumn{4}{p{7.2cm}}{Full mean range 0.129~eV vs.\ 0.425~eV 1-shot; 3.3$\times$ reduction.} & Yes \\
\bottomrule
\end{tabular}
\end{table}

% Table: H1 hypothesis test (15-case)
\begin{table}[t]
\centering
\small
\caption{Explicit test of H$_1$ on the 15-case CMU benchmark using $\Delta E = E_{\mathrm{full}} - E_{\mathrm{1shot}}$. Negative values mean the full agent found a lower energy; hits require $\Delta E < -0.05$~eV. H$_1$ uses only intermetallic surfaces; the four monometallic cases (09, 10, 12, 13) are shown for completeness but excluded from the support rate.}
\label{tab:hypothesis-test}
\begin{tabular}{llrrrr}
\toprule
Case & Family & Gemini $\Delta E$/hit & Grok-4 $\Delta E$/hit & GPT-5.4 $\Delta E$/hit & Claude $\Delta E$/hit \\
\midrule
01 & intermetallic & $-0.001$ / No & $-0.001$ / No & $-0.070$ / Yes & $-0.073$ / Yes \\
02 & intermetallic & $0.000$ / No & $-0.615$ / Yes & $-0.426$ / Yes & $-0.426$ / Yes \\
04 & intermetallic & $-0.402$ / Yes & $-0.597$ / Yes & $-0.012$ / No & $-0.440$ / Yes \\
05 & intermetallic & $-0.015$ / No & $-0.015$ / No & $-0.119$ / Yes & $0.000$ / No \\
09 & monometallic & $-0.173$ / N/A & $0.000$ / N/A & $0.000$ / N/A & $-0.188$ / N/A \\
10 & monometallic & $-0.596$ / N/A & $-0.005$ / N/A & $0.008$ / N/A & $-0.008$ / N/A \\
12 & monometallic & $-0.744$ / N/A & $-0.270$ / N/A & $0.000$ / N/A & $-0.270$ / N/A \\
13 & monometallic & $-0.449$ / N/A & $-0.449$ / N/A & $-0.017$ / N/A & $0.000$ / N/A \\
14 & intermetallic & $0.000$ / No & $-0.372$ / Yes & $-0.160$ / Yes & $-0.126$ / Yes \\
15 & intermetallic & $-0.369$ / Yes & $-0.670$ / Yes & $-0.294$ / Yes & $-0.490$ / Yes \\
16 & intermetallic & $-0.125$ / Yes & $-0.111$ / Yes & $0.391$ / No & $0.000$ / No \\
17 & intermetallic & $-0.449$ / Yes & $0.201$ / No & $-0.549$ / Yes & $-0.549$ / Yes \\
18 & intermetallic & $-0.361$ / Yes & $-0.711$ / Yes & $-0.205$ / Yes & $-0.205$ / Yes \\
19 & intermetallic & $-0.709$ / Yes & $-0.451$ / Yes & $0.000$ / No & $0.000$ / No \\
20 & intermetallic & $-1.048$ / Yes & $-1.050$ / Yes & $0.000$ / No & $-0.000$ / No \\
\midrule
Support rate & --- & 7/11 = 64\% & 8/11 = 73\% & 7/11 = 64\% & 7/11 = 64\% \\
\bottomrule
\end{tabular}
\end{table}
% ===== end inlined paper_tables =====

Tables~\ref{tab:si-ablation-gemini}--\ref{tab:si-ablation-claude} report the full energy matrices for the original locked 15-case subset.
The Friedman test is highly significant on all four backends: Gemini ($p = 1.4 \times 10^{-6}$), Grok-4 ($p = 2.2 \times 10^{-6}$), GPT-5.4 ($p = 2.2 \times 10^{-3}$), and Claude ($p = 1.1 \times 10^{-4}$).
Pairwise Wilcoxon tests confirm that the full agent significantly outperforms the 1-shot baseline after BH correction on three of four backends (Gemini BH~$p = 0.006$, Grok-4 BH~$p = 0.011$, Claude BH~$p = 0.013$); GPT-5.4 is borderline (BH~$p = 0.084$).

\paragraph{The closed-loop gain is large and backend-dependent.}
The mean full-versus-single-shot improvement varies from 0.097~eV (GPT-5.4) to 0.363~eV (Gemini), with Grok-4 at 0.341~eV and Claude at 0.185~eV.
The largest per-case improvements occur on complex surfaces: case~20 shows 1-shot deficits exceeding 1.0~eV on both Gemini and Grok-4, while case~15 consistently degrades under single-shot across all four backends and case~12 degrades on three of four (Gemini, Grok-4, Claude; tied on GPT-5.4).
GPT-5.4 shows the smallest mean improvement, consistent with a stronger one-shot planner that leaves less room for iterative recovery; however, it still gains 0.426~eV on case~02 and 0.549~eV on case~17.

\paragraph{Ablated variants sometimes outperform full.}
On the original locked 15-case subset, the expanded matrix reveals a systematic pattern: ablated variants find lower energies than the full agent on specific case$\times$backend combinations.
Case~18 is the clearest example: \textit{$-$Slip} improves over full by 0.491~eV on GPT-5.4 and Claude, while \textit{$-$Forbid} improves by 0.350~eV on Gemini.
Case~16 shows a reversal where \textit{$-$Slip} degrades by 0.391~eV on Claude but \textit{$-$Forbid}, \textit{$-$Term}, and 1-Shot all match or beat the full variant on GPT-5.4.
Case~17 similarly shows \textit{$-$Slip} outperforming full on Grok-4 by 0.280~eV.
These are not anomalies but a predictable consequence of the guardrail architecture: on surfaces where early attempts happen to find good configurations, the \textit{FORBID} or slip-conditioned constraints can over-prune the remaining search space, steering the agent away from solutions it would otherwise find.
The GPT-5.4 statistics make this concrete: \textit{$-$Forbid} and \textit{$-$Term} both trend toward improving over full (BH~$p = 0.084$, negative rank-biserial $r$), suggesting that on this backend the guardrails impose a small net cost on the expanded case set.

\subsection{Termination mostly controls compute, not chemistry}

The termination signal behaves primarily as a cost-control mechanism across all four backends.
For Grok-4, the mean energy delta between \textit{Full} and \textit{$-$Term} is $-$0.012~eV; disabling termination strictly matches or improves the full agent on 13 of 15 cases (14 of 15 within 0.01~eV).
For Gemini, the same ablation has a mean delta of +0.018~eV, with the largest degradation on case~19 (+0.247~eV).
For GPT-5.4, the effect is marginally beneficial: mean delta $-$0.078~eV, driven by cases~12 and~18 where disabling termination allows the agent to find better solutions.
For Claude, the mean delta is $-$0.031~eV, similarly dominated by case~18 (+0.491~eV improvement).
However, the compute overhead of disabling termination remains substantial: token consumption can increase 2--6$\times$ on individual cases.
Thus, the termination logic should be understood as a guard against planner churn after the search has converged: removing it weakly improves or matches the full agent in many comparisons, but the gains are small, occasional degradations remain, and token use increases on average and sometimes by a large factor.

\subsection{Closed-loop search reduces backend variance}

Table~\ref{tab:cross-llm-summary} summarizes energy agreement across all four backends on the 75 case$\times$variant configurations (15~cases~$\times$~5~variants).
Under the iterative full agent, the mean four-backend range is 0.129~eV, with 7/15~cases agreeing within 0.01~eV.
Under the 1-shot baseline, the mean range rises to 0.425~eV with 0/15~within 0.01~eV---a 3.3$\times$ increase in backend spread.

The tightest convergence is achieved by the \textit{$-$Term} variant (mean range 0.077~eV, 9/15 within 0.01~eV), consistent with the observation that disabling termination gives more cases enough iterations for the four backends to reach similar local minima.
The largest four-backend spread occurs in 1-shot case~20 (1.050~eV), where Gemini and Grok-4 find an energy basin near $-$10.6~eV while GPT-5.4 and Claude converge to $-$11.65~eV.
Under full iterative search, the same case narrows to a range of 0.001~eV---all four backends converge to effectively the same solution.
These data support a concrete systems claim: the closed-loop architecture acts as a backend-robustness mechanism that drives heterogeneous LLM planners toward the same physical solution.

\subsection{Independent validation on the OCD-GMAE dataset}

To test whether the ablation findings generalise beyond the 15 locked CMU cases, we constructed an independent validation set from the Open Catalyst 2020 Dense dataset~\cite{lan2023adsorbml}.
We selected 10 cases spanning 9 distinct adsorbate SMILES (\texttt{[H]}, \texttt{[NH]}, \texttt{[NH2]}, \texttt{N}, \texttt{[N]=N}, \texttt{CO}, \texttt{C(C)O}, \texttt{[CH2][O]}, and \texttt{C([CH2])O}), intermetallic and compound slab systems, and a wide range of adsorption energies ($-$14.55 to $-$1.35~eV).
The same 5$\times$5 ablation matrix was executed successfully on all four LLM backends (Gemini~2.5~Pro, Grok-4, GPT-5.4, Claude~Sonnet~4.6), totaling 200/200 ablation runs. Gemini was accessed via Google Vertex AI direct API.

In parallel, we ran a separate one-shot replication study covering a broader 50-case slice of OCD-GMAE (1 monometallic + 39 intermetallic + 10 compound cases, each evaluated once by each of the same four backends, totalling 200 one-shot runs). This rep50 study is designed to characterise chemical-slip frequency as a function of surface family on this dataset, and is distinct from the 10-case ablation matrix whose energy analysis follows below. Table~\ref{tab:ocd-generalization-summary} summarises the rep50 one-shot outcomes.

\begin{table}[t]
\centering
\scriptsize
\caption{OCD-GMAE rep50 one-shot generalization summary grouped by surface family (50 cases $\times$ 4 backends = 200 one-shot runs, independent of the 10-case ablation matrix). Slip rate is computed as chemical-slip count divided by iteration count for each run.}
\label{tab:ocd-generalization-summary}
\resizebox{\textwidth}{!}{%
\begin{tabular}{lrrrr}
\hline
Surface family & $N$ runs & Success rate & Mean slip rate & Mean iterations  \\
\hline
Monometallic & 4 & 0.0\% & 1.00 & 1.00  \\
Intermetallic & 156 & 89.7\% & 0.78 & 1.00 \\
Compound & 40 & 95.0\% & 0.45 & 1.00  \\
Overall & 200 & 89.0\% & 0.72 & 1.00 \\
\hline
\end{tabular}%
}
\end{table}

% ===== inlined from research/results/ocd_gmae_paper_tables.tex =====
% ============================================================
% OCD-GMAE Independent Validation Paper Tables
% 4-backend analysis: Gemini 2.5 Pro, Grok-4, GPT-5.4, Claude Sonnet 4.6
% Gemini accessed via Vertex AI direct API; remaining three via native endpoints
% Source artifacts:
%   - research/results/ocd_gmae_gemini_ablation_v2/ablation_summary.csv
%   - research/results/ocd_gmae_xai_grok4_ablation_v1/ablation_summary.csv
%   - research/results/ocd_gmae_openai_gpt54_ablation_v1/ablation_summary.csv
%   - research/results/ocd_gmae_anthropic_sonnet46_ablation_v1/ablation_summary.csv
% ============================================================

% Table: OCD-GMAE Ablation Energy Matrix (Gemini 2.5 Pro)
\begin{table}[t]
\centering
\small
\caption{Gemini~2.5~Pro ablation on 10 OCD-GMAE cases drawn from the Open Catalyst 2020 Dense dataset. Lower (more negative) adsorption energy is better; bold marks the best variant per case. Gemini was accessed via Google Vertex AI direct API.}
\label{tab:ocd-ablation-gemini}
\begin{tabular}{lrrrrr}
\toprule
Case & Full & $-$Slip & $-$Forbid & $-$Term & 1-Shot \\
\midrule
003 & \textbf{$-$11.655} & \textbf{$-$11.655} & \textbf{$-$11.655} & \textbf{$-$11.655} & \textbf{$-$11.655} \\
004 & \textbf{$-$10.138} & \textbf{$-$10.138} & \textbf{$-$10.138} & \textbf{$-$10.138} & $-$8.813 \\
005 & $-$10.643 & \textbf{$-$10.643} & \textbf{$-$10.643} & $-$10.607 & \textbf{$-$10.643} \\
012 & \textbf{$-$2.746} & $-$2.694 & \textbf{$-$2.746} & \textbf{$-$2.746} & $-$2.360 \\
013 & \textbf{$-$10.361} & \textbf{$-$10.361} & \textbf{$-$10.361} & \textbf{$-$10.361} & $-$9.077 \\
015 & $-$1.415 & \textbf{$-$1.570} & \textbf{$-$1.570} & $-$1.415 & $-$1.351 \\
016 & \textbf{$-$12.999} & \textbf{$-$12.999} & \textbf{$-$12.999} & \textbf{$-$12.999} & $-$12.921 \\
019 & \textbf{$-$3.504} & \textbf{$-$3.504} & $-$3.338 & $-$3.338 & $-$3.302 \\
023 & \textbf{$-$14.550} & $-$14.256 & \textbf{$-$14.550} & \textbf{$-$14.550} & $-$13.564 \\
024 & \textbf{$-$3.716} & \textbf{$-$3.716} & \textbf{$-$3.716} & \textbf{$-$3.716} & $-$3.377 \\
\bottomrule
\end{tabular}
\end{table}

% Table: OCD-GMAE Ablation Energy Matrix (Grok-4)
\begin{table}[t]
\centering
\small
\caption{Grok-4 ablation on the same 10 OCD-GMAE cases. Bold marks the best variant per case.}
\label{tab:ocd-ablation-grok4}
\begin{tabular}{lrrrrr}
\toprule
Case & Full & $-$Slip & $-$Forbid & $-$Term & 1-Shot \\
\midrule
003 & \textbf{$-$11.659} & $-$11.655 & $-$11.655 & $-$11.655 & $-$10.633 \\
004 & $-$10.138 & $-$10.138 & $-$10.138 & \textbf{$-$11.384} & $-$10.138 \\
005 & $-$10.118 & \textbf{$-$10.643} & \textbf{$-$10.643} & $-$10.607 & $-$9.957 \\
012 & \textbf{$-$2.772} & $-$2.701 & $-$2.746 & $-$2.746 & $-$2.360 \\
013 & \textbf{$-$10.361} & \textbf{$-$10.361} & \textbf{$-$10.361} & \textbf{$-$10.361} & $-$9.078 \\
015 & $-$1.351 & \textbf{$-$1.570} & $-$1.351 & $-$1.351 & $-$1.351 \\
016 & \textbf{$-$12.999} & \textbf{$-$12.999} & \textbf{$-$12.999} & \textbf{$-$12.999} & \textbf{$-$12.999} \\
019 & $-$3.338 & \textbf{$-$3.504} & \textbf{$-$3.504} & \textbf{$-$3.504} & $-$3.338 \\
023 & \textbf{$-$14.550} & $-$14.506 & \textbf{$-$14.550} & \textbf{$-$14.550} & $-$14.256 \\
024 & \textbf{$-$3.716} & $-$3.421 & \textbf{$-$3.716} & \textbf{$-$3.716} & $-$3.377 \\
\bottomrule
\end{tabular}
\end{table}

% Table: OCD-GMAE Ablation Energy Matrix (GPT-5.4)
\begin{table}[t]
\centering
\small
\caption{GPT-5.4 ablation on the same 10 OCD-GMAE cases. Bold marks the best variant per case.}
\label{tab:ocd-ablation-gpt54}
\begin{tabular}{lrrrrr}
\toprule
Case & Full & $-$Slip & $-$Forbid & $-$Term & 1-Shot \\
\midrule
003 & \textbf{$-$11.670} & $-$11.655 & $-$11.655 & $-$11.655 & $-$11.658 \\
004 & \textbf{$-$11.384} & \textbf{$-$11.384} & $-$10.397 & \textbf{$-$11.384} & $-$8.813 \\
005 & \textbf{$-$10.643} & \textbf{$-$10.643} & \textbf{$-$10.643} & \textbf{$-$10.643} & $-$10.011 \\
012 & \textbf{$-$2.769} & $-$2.693 & $-$2.746 & $-$2.694 & $-$2.273 \\
013 & \textbf{$-$10.361} & $-$10.258 & \textbf{$-$10.361} & \textbf{$-$10.361} & $-$8.988 \\
015 & $-$1.551 & $-$1.538 & $-$1.355 & $-$1.559 & \textbf{$-$1.570} \\
016 & \textbf{$-$12.999} & \textbf{$-$12.999} & \textbf{$-$12.999} & \textbf{$-$12.999} & $-$12.921 \\
019 & \textbf{$-$3.504} & $-$3.338 & \textbf{$-$3.504} & \textbf{$-$3.504} & \textbf{$-$3.504} \\
023 & \textbf{$-$14.550} & \textbf{$-$14.550} & \textbf{$-$14.550} & \textbf{$-$14.550} & $-$14.256 \\
024 & \textbf{$-$3.716} & \textbf{$-$3.716} & \textbf{$-$3.716} & \textbf{$-$3.716} & $-$3.377 \\
\bottomrule
\end{tabular}
\end{table}

% Table: OCD-GMAE Ablation Energy Matrix (Claude Sonnet 4.6)
\begin{table}[t]
\centering
\small
\caption{Claude~Sonnet~4.6 ablation on the same 10 OCD-GMAE cases. Bold marks the best variant per case.}
\label{tab:ocd-ablation-claude-ocd}
\begin{tabular}{lrrrrr}
\toprule
Case & Full & $-$Slip & $-$Forbid & $-$Term & 1-Shot \\
\midrule
003 & \textbf{$-$11.655} & \textbf{$-$11.655} & \textbf{$-$11.655} & \textbf{$-$11.655} & $-$6.984 \\
004 & \textbf{$-$11.384} & \textbf{$-$11.384} & \textbf{$-$11.384} & \textbf{$-$11.384} & $-$10.138 \\
005 & \textbf{$-$10.643} & \textbf{$-$10.643} & \textbf{$-$10.643} & \textbf{$-$10.643} & $-$10.209 \\
012 & \textbf{$-$2.746} & $-$2.694 & \textbf{$-$2.746} & \textbf{$-$2.746} & $-$2.273 \\
013 & $-$9.320 & $-$9.670 & $-$9.077 & \textbf{$-$9.673} & $-$9.077 \\
015 & $-$1.351 & \textbf{$-$1.570} & $-$1.351 & $-$1.553 & $-$1.351 \\
016 & \textbf{$-$12.999} & \textbf{$-$12.999} & \textbf{$-$12.999} & \textbf{$-$12.999} & $-$12.914 \\
019 & \textbf{$-$3.504} & $-$3.338 & \textbf{$-$3.504} & \textbf{$-$3.504} & $-$3.257 \\
023 & \textbf{$-$14.550} & \textbf{$-$14.550} & \textbf{$-$14.550} & \textbf{$-$14.550} & $-$14.256 \\
024 & $-$3.418 & \textbf{$-$3.716} & \textbf{$-$3.716} & \textbf{$-$3.716} & $-$3.377 \\
\bottomrule
\end{tabular}
\end{table}

% Table: OCD-GMAE Cross-LLM Robustness (4-backend, Gemini via Vertex AI direct)
\begin{table}[t]
\centering
\caption{Cross-LLM robustness on the OCD-GMAE independent validation set (4~backends: Gemini~2.5~Pro, Grok-4, GPT-5.4, Claude). The four-backend range is max--min energy for the same case and variant.}
\label{tab:ocd-cross-llm}
\begin{tabular}{lrrr}
\toprule
Variant & Mean 4-backend range (eV) & $\leq$0.01\,eV & Largest range (eV) \\
\midrule
Full & 0.352 & 2/10 & 1.247 \\
$-$Slip & 0.273 & 4/10 & 1.247 \\
$-$Forbid & 0.292 & 6/10 & 1.284 \\
$-$Term & 0.240 & 4/10 & 1.247 \\
1-Shot & 0.810 & 1/10 & 4.674 \\
\midrule
Overall & 0.393 & 17/50 & 4.674 \\
\bottomrule
\end{tabular}
\end{table}
% ===== end inlined ocd_gmae_paper_tables =====

\paragraph{Closed-loop improvement is larger and more consistent.}
The mean full-versus-single-shot improvement across the 10 cases is 0.542~eV (4-backend average), substantially larger than the 0.247~eV observed on the CMU benchmark.
Nine of ten cases show improvement exceeding 0.05~eV; only case~015 is effectively neutral (mean single-shot minus full $\Delta = +0.011$~eV across backends).
The largest gains occur on cases with multiple competitive binding sites: case~004 (+1.325~eV on Gemini), case~013 (+1.284~eV on Gemini), and case~023 (+0.985~eV on Gemini).

\paragraph{Statistical significance is achieved.}
With 10 cases, the Friedman test reaches significance on all four backends: Gemini ($p = 3.4 \times 10^{-4}$), Grok-4 ($p = 2.0 \times 10^{-3}$), GPT-5.4 ($p = 6.9 \times 10^{-3}$), Claude ($p = 8.2 \times 10^{-5}$).
For GPT-5.4, the pairwise Wilcoxon test for $-$Slip versus Full is individually significant ($p = 0.016$), indicating that slip-conditioned feedback contributes measurably to search quality on this dataset.

\paragraph{Backend convergence pattern.}
Table~\ref{tab:ocd-cross-llm} reports the four-backend cross-LLM robustness analysis (Gemini accessed via Vertex AI direct API, the remaining three via their native endpoints).
The $-$Term variant achieves the tightest cross-backend agreement (mean range 0.240~eV), while single-shot shows the widest spread (0.810~eV)---a 3.4$\times$ ratio.
The overall pattern mirrors the CMU result but with higher absolute ranges, reflecting the greater chemical diversity and more complex potential-energy surfaces in the OCD-GMAE set.
The Full variant shows a mean range of 0.352~eV, driven primarily by case~004, where Grok-4 converges to only $-$10.138~eV while GPT-5.4 and Claude reach $-$11.384~eV (range 1.247~eV), and by case~013, where Claude finds a distinct local minimum at $-$9.320~eV that the other three backends miss ($-$10.361~eV).

\paragraph{Ablated variants can outperform full.}
Several cases show ablated variants finding lower energies than Full: case~004 under Grok-4 $-$Term ($-$11.384 vs.\ $-$10.138~eV, $\Delta = -1.247$~eV), case~005 under Grok-4 $-$Slip ($-$10.643 vs.\ $-$10.118~eV), and case~024 under Claude $-$Slip ($-$3.716 vs.\ $-$3.418~eV).
These cases support the CMU observation that the agentic guardrails are not universally beneficial: on certain surfaces, the FORBID or slip-conditioned constraints over-prune the search space.

\subsection{Comparison with Adsorb-Agent under identical physics}

To isolate the effect of search strategy from the choice of force field, we forked the open-source Adsorb-Agent codebase~\cite{adsorbagent} and replaced its EquiformerV2 relaxation backend with the same MACE-MP-0 (small, CPU, float32) calculator used for the primary benchmark and matched comparisons.
All remaining protocol parameters were locked to AdsMind's settings: BFGS optimizer, $f_\mathrm{max} = 0.10$~eV/\AA, 200 steps, and identical \textsc{FixAtoms} constraints (bottom third of slab atoms fixed).
We ran the modified Adsorb-Agent with GPT-5.4 and report the 15-case overlap with the locked CMU ablation subset.
We separately ran a single-configuration Adsorb-Agent stress control on CMU20 to test whether Adsorb-Agent's energy depth comes from one-shot chemical site choice alone or from its breadth of enumerated configurations.

\begin{table}[t]
\centering
\scriptsize
\caption{Matched-physics comparison between AdsMind Full and Adsorb-Agent on the 15 locked CMU benchmark cases using GPT-5.4 and MACE-MP-0 small (CPU, float32). Energies are adsorption energies in eV; NULL denotes no valid Adsorb-Agent trajectory.}
\label{tab:cmu-gpt54-adsorbagent-comparison}
\resizebox{\textwidth}{!}{%
\begin{tabular}{llllrrrr}
\hline
Case & Surface & Adsorbate & AdsMind success & AdsMind energy & AdsMind iters & Adsorb-Agent energy & Adsorb-Agent configs \\
\hline
01 & Mo3Pd & H & Yes & $-$3.627 & 4 & $-$3.879 & 29 \\
02 & Mo3Pd & NNH & Yes & $-$4.769 & 5 & $-$5.097 & 2 \\
04 & CuPd3 & NNH & Yes & $-$2.255 & 4 & $-$2.915 & 34 \\
05 & Cu3Ag & H & Yes & $-$2.962 & 4 & NULL & 0 \\
09 & Pt & OH & Yes & $-$1.991 & 3 & $-$1.997 & 25 \\
10 & Pt & OH & Yes & $-$2.710 & 2 & $-$2.741 & 22 \\
12 & Au & OH & Yes & $-$2.081 & 1 & $-$2.245 & 23 \\
13 & Ag & OH & Yes & $-$2.825 & 3 & $-$2.857 & 25 \\
14 & CoPt & OH & Yes & $-$3.591 & 5 & $-$4.125 & 8 \\
15 & Cu6Ga2 & CH$_2$CH$_2$OH & Yes & $-$3.258 & 4 & NULL & 0 \\
16 & Au2Hf & CH$_2$CH$_2$OH & Yes & $-$4.272 & 5 & $-$5.700 & 20 \\
17 & Rh2Ti2 & OCHCH$_3$ & Yes & $-$3.241 & 5 & $-$4.476 & 31 \\
18 & Al3Zr & OCHCH$_3$ & Yes & $-$2.724 & 3 & NULL & 0 \\
19 & Hf2Zn6 & OCHCH$_3$ & Yes & $-$4.013 & 5 & $-$4.707 & 18 \\
20 & Bi2Ti6 & ONN(CH$_3$)$_2$ & Yes & $-$11.652 & 5 & $-$12.285 & 16 \\
\hline
\end{tabular}%
}
\end{table}

\paragraph{Reliability gap.}
AdsMind succeeds on all 15 overlapping cases with the GPT-5.4 full variant (100\%).
Adsorb-Agent succeeds on 12 of these 15 cases (80\%) and fails on cases~05, 15, and 18 because no valid trajectory files are produced.
This failure mode is consistent with a single-pass architecture: if the initial enumerated candidate set is empty or invalid, the system has no feedback loop that can redirect the planner, whereas AdsMind's closed-loop diagnostics can continue after failed or slipped attempts.

\paragraph{Energy depth under matched physics.}
Among the 12 cases where both systems produce valid energies, Adsorb-Agent finds a lower adsorption energy on all 12.
With $\Delta E = E_{\mathrm{AdsMind}} - E_{\mathrm{AdsorbAgent}}$, the mean difference is $+0.500$~eV in favor of Adsorb-Agent (Wilcoxon $p = 4.88 \times 10^{-4}$, $n = 12$).
The largest gaps occur on cases with larger configuration spaces: case~16 ($\Delta = 1.428$~eV), case~17 ($\Delta = 1.235$~eV), and case~19 ($\Delta = 0.695$~eV).
On simpler OH cases, the two systems can be nearly indistinguishable, e.g., case~09 differs by only 0.006~eV.

\paragraph{The gap is explained by sampling volume.}
Adsorb-Agent relaxes 2--34 initial configurations on successful cases (median~22.5, mean~21.1), evaluating many candidate placements in parallel.
AdsMind relaxes one configuration per iteration over 1--5 iterations (mean~3.9) in the GPT-5.4 full run.
This $\sim$5.5$\times$ relaxation-count difference provides a mechanistic explanation for why Adsorb-Agent reaches deeper energy basins when it succeeds: on surfaces with many local minima, broader enumeration can find minima that a single-trajectory closed loop may never enter.
The single-configuration control supports this interpretation: when Adsorb-Agent is forced to select only one configuration, it produces no selected configuration on 6/20 CMU cases and is worse than its original multi-configuration run on all 12 cases with comparable derived adsorption energies (mean single-minus-multi gap $+1.109$~eV).
We therefore interpret Adsorb-Agent's energy advantage as primarily a breadth effect, not as evidence that a single open-loop LLM placement is intrinsically deeper than closed-loop refinement.

\paragraph{Cost structure.}
The two systems have inverted cost profiles.
Adsorb-Agent consumes roughly 5{,}100 LLM tokens on successful cases but about 21 MACE relaxations per case; AdsMind consumes more LLM tokens through iterative reasoning but only about four MACE relaxations per case.
The current matched-physics comparison therefore should not be read as AdsMind winning on raw energy depth.
It supports a reliability--cost trade-off: AdsMind is more reliable and uses far fewer relaxations, while Adsorb-Agent is deeper when its one-pass enumeration succeeds.

\paragraph{Implications.}
AdsMind's closed-loop design achieves successful Full-variant convergence across the completed CMU20 benchmark, but its iterative single-configuration refinement under-samples the configuration space on complex surfaces.
Adsorb-Agent's multi-configuration enumeration explores more of the energy landscape when the candidate set is valid, but fails when the one-pass setup produces no valid trajectories.
These findings motivate a concrete architectural extension: coupling AdsMind's closed-loop error recovery with broader initial site enumeration and parallel relaxation.

\subsection{Non-LLM and cross-backend control baselines}

To isolate the contribution of the LLM planner from that of the underlying MACE relaxation pipeline, we implemented two non-LLM controls using the same MACE-matched physics: a \textit{random placement} baseline (N=20 random adsorbate poses per case, each relaxed with the standard protocol) and a \textit{heuristic enumeration} baseline (all high-symmetry sites identified by Autoadsorbate, median 53 sites per case).
We include the GPT-5.4 Adsorb-Agent run above and in the SI baseline table to keep the LLM backend matched where possible.

\paragraph{Random placement is competitive with AdsMind Full when given more samples.}
After adding the five previously missing random-control cases, the CMU20 comparison gives a mean energy difference of $-0.232$~eV relative to AdsMind Full (best-across-backends); random wins 6 cases, AdsMind wins 4, and 10 tie within 0.01~eV.
The two largest random advantages remain case~16 ($-6.74$ vs.\ $-4.66$ eV, $\Delta = -2.07$~eV) and case~20 ($-13.42$ vs.\ $-11.65$ eV, $\Delta = -1.77$~eV), while AdsMind's largest advantage is case~19 ($+0.477$~eV in random-minus-AdsMind energy).
Critically, the random baseline uses $5\times$ more MACE relaxations than a single AdsMind Full run on the same case (20 vs.\ $\sim$4); the best-across-backends comparison above is a reporting view for robustness, not a claim that the four-backend aggregate itself costs only $\sim$4 relaxations.
In the CMU20 case-level best-of-backend view, AdsMind's \textit{1-Shot} variant, which allots the same single MACE relaxation as a naive search would, never beats AdsMind Full and underperforms random on several cases.
This supports the interpretation that the LLM planner's single-pass placement is not inherently better than broader random sampling; the value of the agent lies in \textit{iterative refinement}, reliability, and interpretable error recovery rather than raw deepest-basin discovery.
Because the current random summary records successful MACE energy calculations but does not yet apply the same anomaly filtering as the agent summaries, random winners should be interpreted as protocol-level energy comparisons until their top structures are manually or programmatically filtered.

\paragraph{Heuristic enumeration beats AdsMind Full on 6 of 15 cases.}
Systematically relaxing every Autoadsorbate-enumerated site (25--98 sites per case, no LLM) wins on 6 cases, loses on 1, ties on 8.
On case~19 (Hf$_2$Zn$_6$ + OCHCH$_3$), heuristic finds $-4.70$~eV versus AdsMind Full's $-4.04$~eV---a 0.65~eV gap comparable to Adsorb-Agent's.
The heuristic baseline uses 6--25$\times$ more MACE compute than AdsMind Full.
This result positions AdsMind's LLM planner as a \textit{compute-efficient approximation} of exhaustive site enumeration: AdsMind trails the heuristic optimum by 0.19~eV on average while using 1/6 to 1/25 the number of relaxations.

\paragraph{Adsorb-Agent LLM choice changes failures but not the trade-off.}
The current GPT-5.4 Adsorb-Agent run succeeds on 12 of 15 locked CMU cases, while the historical GPT-4o run succeeds on 11 of 15.
Failure cases differ between LLMs (GPT-4o fails on 01, 09, 15, 20; GPT-5.4 fails on 05, 15, 18), suggesting that single-pass LLM planning retains a $\sim$20--27\% failure rate even as the backend changes.
AdsMind's closed-loop architecture shows a 0\% failure rate across all four backends on the same locked subset.

\paragraph{Iteration convergence.}
Tracing the per-iteration running-best energy across all 15 cases and 4 backends reveals that iteration 2 captures 58--73\% of the eventual full-run improvement over iteration~1, and iteration~3 captures 74--90\%.
The remaining iterations contribute diminishing returns.
This rapid convergence is the mechanistic basis for AdsMind's compute efficiency: most of the closed-loop value is extracted in the first 2--3 iterations.

\paragraph{MACE-large sensitivity.}
Repeating the GPT-5.4 Full variant on all 20 CMU cases with MACE-large (GPU, float64, dispersion) shows that force-field fidelity can shift absolute adsorption energies substantially.
All 20 CMU rows produce retained summary results, but the mean absolute shift relative to the MACE-small GPT-5.4 run is 1.255~eV (0.834~eV if the anomalous case~20 is excluded), and no case remains within 0.01~eV.
Fifteen cases become more stable under MACE-large and five become less stable.
Cases~06, 16, 17, and 20 show dissociation events during the MACE-large runs; case~20 is the most important sensitivity case because a retained non-dissociated best structure appears at $-2.402$~eV while a later dissociated analysis reaches a much lower unretained energy.
On a separate 10-case OCD-rep50 representative set, MACE-large succeeds on 9/10 cases; case~043 fails naturally with chemical slip and dissociation, not because of an API or calculator error.
These shifts do not change the qualitative systems conclusion from the MACE-small benchmark: primary comparisons must remain within the locked MACE-small protocol, and AdsMind's strength is reliable closed-loop recovery with a small relaxation budget rather than force-field-independent deepest-basin discovery.

\paragraph{Reframing AdsMind's contribution.}
Taken together, these controls reveal that AdsMind's primary contribution is not deepest-basin discovery.
Random sampling with 20 relaxations is competitive with AdsMind Full and reaches a lower mean energy; heuristic enumeration with $\sim$50 sites beats it on 40\% of cases; Adsorb-Agent with $\sim$21 relaxations reaches deeper basins whenever the GPT-5.4 one-pass enumeration succeeds.
In the GPT-5.4 matched-physics run, Adsorb-Agent is lower in energy on all 12 successful comparable cases, but fails on the remaining three.
AdsMind's distinctive value is the combination of: (i) 100\% success for the Full variant on both completed ablation matrices (80/80 CMU20 and 40/40 OCD-GMAE), with 320/320 successful CMU20 iterative-variant runs versus 73/80 for the CMU20 1-Shot control, 73--80\% for Adsorb-Agent on the locked CMU overlap, and 14/20 selected configurations in the Adsorb-Agent single-configuration stress test; (ii) compute efficiency per backend---remaining competitive with random at $5\times$ less MACE work per run and approaching heuristic at $6$--$25\times$ less; (iii) chemical interpretability via Slip detection, FORBID logic, and backend-convergent reasoning traces absent from brute-force methods; and (iv) iterative error recovery, which is a prerequisite for autonomous chemistry workflows where silent failures are unacceptable.
A natural extension couples AdsMind's closed-loop error-recovery with broader initial sampling (multi-site LLM placement with parallel relaxation), which would combine the depth of enumeration with the reliability and interpretability of the closed-loop agent.
